\documentclass[../../main.tex]{subfiles}% 注意这里的文件路径不能用 ./main.tex ,否则用latexmk编译子文件会报错
\graphicspath{{\subfix{./image/}}} % 指定图片目录,后续可以直接使用图片文件名
% 注意这里的文件路径不能用 ../../image/ ,否则用latexmk编译子文件会报错

% 例如:
% \begin{figure}[H]
% \centering
% \includegraphics[scale=0.3]{图.png}
% \caption{}
% \label{figure:图}
% \end{figure}
% 注意:上述\label{}一定要放在\caption{}之后,否则引用图片序号会只会显示??.

\begin{document}

\section{代数学基本定理}

设首一$n$次多项式
\begin{align}
f(x)=x^n+a_{n-1}x^{n-1}+\dots+a_1x+a_0=(x-\alpha_1)(x-\alpha_2)\cdots(x-\alpha_n),\label{eq:poly-root}
\end{align}
展开后可得根的初等对称多项式与系数的关系：
\begin{align*}
\begin{cases}
\alpha_1+\alpha_2+\dots+\alpha_n=-a_{n-1},\\
\displaystyle\sum\limits_{i<j}\alpha_i\alpha_j=a_{n-2},\\
\quad\cdots\cdots\\
\displaystyle\sum\limits_{i_1<i_2<\dots<i_k}\alpha_{i_1}\alpha_{i_2}\cdots\alpha_{i_k}=(-1)^k a_{n-k},\\
\quad\cdots\cdots\\
\alpha_1\alpha_2\cdots\alpha_n=(-1)^n a_0.
\end{cases}
\end{align*}
由对称多项式基本定理，关于$f(x)$的根的任意对称多项式均可表示为$f(x)$系数的多项式。

\begin{lemma}
奇数次实多项式有实数根。
\end{lemma}

\begin{proof}
设实多项式$f(x)=x^n+a_{n-1}x^{n-1}+\dots+a_1x+a_0$，其中$n$为奇数。
记$A=\max\{|a_0|,|a_1|,\dots,|a_{n-1}|\}$。
当$|x|>1$时，
\begin{align*}
|a_{n-1}x^{n-1}+\dots+a_1x+a_0|&\leqslant|a_{n-1}||x|^{n-1}+\dots+|a_1||x|+|a_0|\\
&\leqslant A\cdot\frac{|x|^n-1}{|x|-1}\\
&<A\cdot\frac{|x|^n}{|x|-1}.
\end{align*}
当$|x|\geqslant A+1$时，$|x|^n>|a_{n-1}x^{n-1}+\dots+a_1x+a_0|$，因此$f(|x|)>0$，$f(-|x|)<0$。
$f(x)$为$\mathbb{R}$上连续函数，由介值定理，存在$\alpha\in\mathbb{R}$使得$f(\alpha)=0$.

\end{proof}

\begin{lemma}
次数大于零的实多项式有复数根。
\end{lemma}

\begin{proof}
对实多项式$f(x)=x^n+a_{n-1}x^{n-1}+\dots+a_1x+a_0$，设$\deg f=2^k q$，其中$q$为奇数，对$k$作数学归纳法。

当$k=0$时，$n=q$为奇数，由上一引理结论成立。
假设$k-1$时结论成立，现证$n=2^k q$时结论成立。

将$\mathbb{R}$扩域至包含$f(x)$全部根$\alpha_1,\alpha_2,\dots,\alpha_n$的域$K$。对任意实数$c$，令
\begin{align}
\beta_{ij}=\alpha_i\alpha_j+c(\alpha_i+\alpha_j),\quad 1\leqslant i<j\leqslant n,\label{eq:beta-def}
\end{align}
构造多项式
\begin{align}
g(x)=\prod\limits_{1\leqslant i<j\leqslant n}(x-\beta_{ij}).\label{eq:g-poly}
\end{align}
计算其次数：
$$\deg g=\frac12 n(n-1)=2^{k-1}q(2^k q-1)=2^{k-1}q',$$
其中$q'$为奇数。

$g(x)$的系数为$\beta_{ij}$的实对称多项式，而$\beta_{ij}$是$\alpha_i,\alpha_j$的实多项式，故$g(x)$的系数是$\alpha_1,\dots,\alpha_n$的实对称多项式。由对称多项式基本定理，$g(x)$的系数可表示为$f(x)$系数的实多项式，故$g(x)$为实多项式。

由归纳假设，$g(x)$存在复数根，即存在实数$c$与下标$\{i,j\}$，使得$\alpha_i\alpha_j+c(\alpha_i+\alpha_j)$为复数。取两个不同实数$c_1,c_2$，使得
\begin{align}
\begin{cases}
\alpha_i\alpha_j+c_1(\alpha_i+\alpha_j)=a,\\
\alpha_i\alpha_j+c_2(\alpha_i+\alpha_j)=b,
\end{cases}\label{eq:sys-ab}
\end{align}
其中$a,b$均为复数。解得
$$\alpha_i+\alpha_j=\frac{a-b}{c_1-c_2},\quad \alpha_i\alpha_j=a-c_1(\alpha_i+\alpha_j),$$
二者均为复数。因此二次多项式$x^2-(\alpha_i+\alpha_j)x+\alpha_i\alpha_j$的根$\alpha_i,\alpha_j$均为复数，即$f(x)$有复数根.

\end{proof}

\begin{theorem}
次数大于零的复系数一元多项式有复数根。
\end{theorem}

\begin{proof}
设复系数多项式$f(x)=x^n+a_{n-1}x^{n-1}+\dots+a_1x+a_0$，记$\overline{f}(x)=x^n+\overline{a}_{n-1}x^{n-1}+\dots+\overline{a}_1x+\overline{a}_0$，其中$\overline{a}_i$为$a_i$的共轭复数。
构造多项式
\begin{align}
F(x)=f(x)\overline{f}(x)=x^{2n}+b_{2n-1}x^{2n-1}+\dots+b_1x+b_0.\label{eq:F-poly}
\end{align}
其系数满足
$$b_k=\sum\limits_{i+j=k}a_i\overline{a}_j=\sum\limits_{i+j=k}\overline{a}_i a_j=\overline{b}_k,\quad 0\leqslant k\leqslant 2n-1,$$
故$F(x)$为实多项式。

由引理，$F(x)$存在复数根$\beta$，即$F(\beta)=f(\beta)\overline{f}(\beta)=0$。
若$f(\beta)=0$，则$\beta$为$f(x)$的根；若$f(\beta)\neq0$，则$\overline{f}(\beta)=0$，计算得
\begin{align*}
0&=\overline{f}(\beta)=\beta^n+\overline{a}_{n-1}\beta^{n-1}+\dots+\overline{a}_1\beta+\overline{a}_0\\
&=\overline{\overline{\beta}^n+a_{n-1}\overline{\beta}^{n-1}+\dots+a_1\overline{\beta}+a_0}=\overline{f(\overline{\beta})},
\end{align*}
故$f(\overline{\beta})=0$，即$\overline{\beta}$为$f(x)$的复数根.

\end{proof}



\end{document}